%%fssp_ncor_specrs

L'algorithme et les formules \'etant complexes
seul le principe de calcul est pr\'esent\'e ici.\\
Selon Brenguier-89 \cite{Brenguier-89}, 
le nombre de gouttes mesur\'ees par un FSSP
s'exprime comme une fonction~: \\ 
\begin{eqnarray}
N & = & f_n ( N_c, \tau, \tau_1, \tau_2, \tau_r ) \label{eqn:crt2}
\end{eqnarray}
o\`u~:\\
$N$ : nombre total de gouttes compt\'ees \\
$N_c$ : nombre total de gouttes corrig\'e des
co\"{\i}ncidences \\
$\tau $ : temps de transit d'une particule \\
\\
On pose $\tau = \tau_{spec}$, et $N=Reset$.
Ainsi, cette \'equation ne
poss\`ede plus qu'une inconnue $N_c$. 
Elle est r\'esolue par it\'erations successives.   
(\ref{eqn:crt2}) donne $N_{rs}$.\\
Si $Reset\times f_e \le Crt_0$ alors aucune correction
est calcul\'ee et $N_{ac}=0$ et $N_{rs}=0$.
